\chapter{Lesson 9: Health and the Hospital}

This lesson covers describing symptoms at a clinic or pharmacy and understanding the
advice you get. The grammar is the conditional \ko{\til (으)면} (if, when), obligation
\ko{\til 아/어야 하다/되다} (must), the related permission and prohibition patterns, the
negative command \ko{\til 지 마세요}, and the \ko{ㅡ}-irregular verbs, which include the
key verb of the lesson, \ko{아프다} (to hurt, be sick).

\section{Vocabulary}

\subsection{The body}

\begin{center}
\begin{tikzpicture}[every node/.style={font=\scriptsize}, lab/.style={draw=none, fill=white, inner sep=1pt}]
  % stick figure
  \draw[thick] (0,3.6) circle (0.5);
  \draw[thick] (0,3.1) -- (0,1.2);
  \draw[thick] (0,2.7) -- (-1.1,1.7) -- (-1.3,1.2);
  \draw[thick] (0,2.7) -- (1.1,1.7) -- (1.3,1.2);
  \draw[thick] (0,1.2) -- (-0.6,0) -- (-0.8,-0.1);
  \draw[thick] (0,1.2) -- (0.6,0) -- (0.8,-0.1);
  \fill (-0.18,3.7) circle (1pt); \fill (0.18,3.7) circle (1pt);
  % left labels
  \draw[gray] (-0.45,3.9) -- (-2.2,4.2) node[left] {\ko{머리} meori (head)};
  \draw[gray] (-0.2,3.7) -- (-2.2,3.6) node[left] {\ko{눈} nun (eye)};
  \draw[gray] (0,2.95) -- (-2.2,3.0) node[left] {\ko{목} mok (neck, throat)};
  \draw[gray] (-0.7,2.35) -- (-2.2,2.4) node[left] {\ko{팔} pal (arm)};
  \draw[gray] (-1.3,1.2) -- (-2.2,1.2) node[left] {\ko{손} son (hand)};
  \draw[gray] (-0.4,0.45) -- (-2.2,0.4) node[left] {\ko{다리} dari (leg)};
  % right labels
  \draw[gray] (0.0,3.45) -- (2.2,4.0) node[right] {\ko{코} ko (nose), \ko{입} ip (mouth)};
  \draw[gray] (0.5,3.6) -- (2.2,3.4) node[right] {\ko{귀} gwi (ear)};
  \draw[gray] (0.4,2.7) -- (2.2,2.8) node[right] {\ko{어깨} eokkae (shoulder)};
  \draw[gray] (0,1.9) -- (2.2,2.0) node[right] {\ko{배} bae (stomach)};
  \draw[gray] (0,1.4) -- (2.2,1.3) node[right] {\ko{허리} heori (lower back)};
  \draw[gray] (0.8,-0.1) -- (2.2,0.2) node[right] {\ko{발} bal (foot)};
\end{tikzpicture}
\end{center}

\subsection{Symptoms and the clinic}

\begin{table}[H]
\centering
\begin{tabular}{lll}
\toprule
Korean & RR & English \\
\midrule
\ko{아프다} & apeuda & to hurt; to be sick \\
\ko{감기에 걸리다} & gamgie geollida & to catch a cold \\
\ko{열이 나다} & yeori nada & to have a fever \\
\ko{기침을 하다} & gichimeul hada & to cough \\
\ko{콧물이 나다} & konmuri nada & to have a runny nose \\
\ko{목이 붓다} & mogi butda & to have a swollen (sore) throat \\
\ko{배탈이 나다} & baetari nada & to have an upset stomach \\
\ko{다치다} & dachida & to get hurt, injured \\
\ko{낫다} & natda & to get better, recover \\
\ko{병원} & byeongwon & hospital, clinic \\
\ko{약국} & yakguk & pharmacy \\
\ko{약} & yak & medicine \\
\ko{약을 먹다} & yageul meokda & to take medicine \\
\ko{주사를 맞다} & jusareul matda & to get a shot \\
\ko{의사 / 간호사} & uisa / ganhosa & doctor / nurse \\
\ko{푹 쉬다} & puk swida & to rest well \\
\ko{조심하다} & josimhada & to be careful \\
\bottomrule
\end{tabular}
\end{table}

\begin{remark}
``My head hurts'' is \ko{머리가 아파요} --- the body part is the \emph{subject}. To say
where it hurts on a person, use the double-subject pattern: \ko{저는 머리가 아파요} (as
for me, my head hurts). \ko{붓다} and \ko{낫다} are \ko{ㅅ}-irregular (\ko{부었어요},
\ko{나았어요}); see the next chapter.
\end{remark}

\section{Grammar 1: \ko{ㅡ}-irregular verbs}

When a stem ends in the vowel \ko{ㅡ} and the ending begins with \ko{아/어}, the
\ko{ㅡ} drops. Which ending you get depends on the vowel of the \emph{preceding}
syllable; one-syllable stems take \ko{어}.

\begin{table}[H]
\centering
\begin{tabular}{lllll}
\toprule
Dictionary & Preceding vowel & \ko{\til 아/어요} & Past & Meaning \\
\midrule
\ko{아프다} & \ko{ㅏ} & \ko{아파요} & \ko{아팠어요} & hurt, be sick \\
\ko{바쁘다} & \ko{ㅏ} & \ko{바빠요} & \ko{바빴어요} & be busy \\
\ko{배고프다} & \ko{ㅗ} & \ko{배고파요} & \ko{배고팠어요} & be hungry \\
\ko{예쁘다} & \ko{ㅖ} & \ko{예뻐요} & \ko{예뻤어요} & be pretty \\
\ko{슬프다} & \ko{ㅡ} & \ko{슬퍼요} & \ko{슬펐어요} & be sad \\
\ko{크다} & (none) & \ko{커요} & \ko{컸어요} & be big \\
\ko{쓰다} & (none) & \ko{써요} & \ko{썼어요} & write; use; wear (hat) \\
\bottomrule
\end{tabular}
\end{table}

\section{Grammar 2: \ko{\til (으)면} (if, when)}

\ko{\til (으)면} attaches to the stem: \ko{면} after a vowel or \ko{ㄹ}, \ko{으면} after
other consonants. It expresses a condition (if) or a general ``whenever'' relationship.
The result clause can be a statement, question, command or suggestion.

\begin{example}[Conditions]
A doctor's advice and general truths:
\begin{itemize}
  \item \ko{열이 나면 이 약을 드세요.} Yeori namyeon i yageul deuseyo. (If you have a
        fever, take this medicine.)
  \item \ko{감기에 걸리면 물을 많이 마시세요.} Gamgie geollimyeon mureul mani masiseyo.
        (When you catch a cold, drink a lot of water.)
  \item \ko{계속 아프면 다시 오세요.} Gyesok apeumyeon dasi oseyo. (If it keeps hurting,
        come back.)
  \item \ko{시간이 있으면 같이 병원에 갈까요?} Sigani isseumyeon gachi byeongwone
        galkkayo? (If you have time, shall we go to the clinic together?)
  \item \ko{학교가 멀면 버스를 타세요.} Hakgyoga meolmyeon beoseureul taseyo. (If school
        is far, take the bus.) --- \ko{ㄹ}-stem + \ko{면}.
  \item \ko{날씨가 추우면 감기에 잘 걸려요.} Nalssiga chuumyeon gamgie jal geollyeoyo.
        (When it's cold, I catch colds easily.) --- \ko{ㅂ}-irregular: \ko{춥다} $\to$
        \ko{추우면}.
\end{itemize}
\end{example}

\section{Grammar 3: \ko{\til 아/어야 하다/되다} (must, have to)}

Attach \ko{\til 아/어야} (vowel harmony, like \ko{\til 아/어요}) and then \ko{하다} or
\ko{되다}. The two are interchangeable; \ko{되다} is a little more colloquial.

\begin{example}[Obligation]
\ko{\til 아/어야 해요} / \ko{\til 아/어야 돼요}:
\begin{itemize}
  \item \ko{하루에 세 번 약을 먹어야 해요.} Harue se beon yageul meogeoya haeyo. (You must
        take the medicine three times a day.)
  \item \ko{식사 후에 드셔야 돼요.} Siksa hue deusyeoya dwaeyo. (You have to take it after
        meals.) --- honorific \ko{드시다}.
  \item \ko{오늘은 집에서 푹 쉬어야 해요.} Oneureun jibeseo puk swieoya haeyo. (You
        should rest well at home today.)
  \item \ko{어제 병원에 가야 했어요.} Eoje byeongwone gaya haesseoyo. (Yesterday I had to
        go to the hospital.) --- past tense goes on \ko{하다}.
\end{itemize}
\end{example}

\section{Grammar 4: Permission, prohibition, and ``don't''}

Three more patterns complete the set of what you must, may, need not, and must not do.
\begin{itemize}
  \item \textbf{May}: \ko{\til 아/어도 되다}: \ko{샤워해도 돼요?} Syawohaedo dwaeyo? (May
        I take a shower?)
  \item \textbf{Must not}: \ko{\til (으)면 안 되다}, literally ``if you do it, it's no
        good'': \ko{술을 마시면 안 돼요.} Sureul masimyeon an dwaeyo. (You must not drink
        alcohol.)
  \item \textbf{Need not}: \ko{\til 지 않아도 되다} or \ko{안 \til 아/어도 되다}: \ko{내일은
        안 와도 돼요.} Naeireun an wado dwaeyo. (You don't have to come tomorrow.)
  \item \textbf{Don't} (negative command): \ko{\til 지 마세요}: \ko{무리하지 마세요.}
        Murihaji maseyo. (Don't overdo it.) \ko{걱정하지 마세요.} Geokjeonghaji maseyo.
        (Don't worry.)
\end{itemize}

\begin{center}
\begin{tikzpicture}[every node/.style={font=\small, align=center}]
  \draw[thick] (0,0) rectangle (8,4);
  \draw (4,0) -- (4,4); \draw (0,2) -- (8,2);
  \node[font=\scriptsize] at (2,4.3) {do it};
  \node[font=\scriptsize] at (6,4.3) {don't do it};
  \node[font=\scriptsize, rotate=90] at (-0.3,3) {required};
  \node[font=\scriptsize, rotate=90] at (-0.3,1) {optional};
  \fill[red!10] (0.05,2.05) rectangle (3.95,3.95);
  \fill[green!10] (0.05,0.05) rectangle (3.95,1.95);
  \fill[red!10] (4.05,2.05) rectangle (7.95,3.95);
  \fill[green!10] (4.05,0.05) rectangle (7.95,1.95);
  \node at (2,3) {\ko{\til 아/어야 해요}\\ \scriptsize must};
  \node at (2,1) {\ko{\til 아/어도 돼요}\\ \scriptsize may};
  \node at (6,3) {\ko{\til (으)면 안 돼요}\\ \ko{\til 지 마세요}\\ \scriptsize must not / don't};
  \node at (6,1) {\ko{\til 지 않아도 돼요}\\ \scriptsize don't have to};
\end{tikzpicture}
\end{center}

\section{At the clinic: a model conversation}

\begin{quote}
\ko{의사: 어디가 아프세요?}\\
\ko{환자: 어제부터 머리가 아프고 열이 나요. 기침도 많이 해요.}\\
\ko{의사: 목도 아프세요? 아 하세요. \ldots 목이 많이 부었네요. 감기예요.}\\
\ko{환자: 주사를 맞아야 돼요?}\\
\ko{의사: 아니요, 주사는 안 맞아도 돼요. 이 약을 하루에 세 번 식사 후에 드세요.}\\
\ko{환자: 학교에 가도 돼요?}\\
\ko{의사: 열이 있으면 학교에 가지 마세요. 집에서 푹 쉬어야 해요. 그리고 찬 음료수를 마시면
안 돼요.}\\
\ko{환자: 네, 감사합니다.}\\
\ko{의사: 빨리 나으세요.}
\end{quote}

\noindent Doctor: Where does it hurt? / Patient: Since yesterday I've had a headache and
a fever. I'm coughing a lot too. / Doctor: Does your throat hurt? Say ``ah''. \ldots
Your throat is quite swollen. It's a cold. / Patient: Do I need a shot? / Doctor: No,
you don't need a shot. Take this medicine three times a day after meals. / Patient: Can
I go to school? / Doctor: If you have a fever, don't go to school. You must rest well at
home. And you mustn't drink cold drinks. / Patient: OK, thank you. / Doctor: Get well
soon.

\begin{aside}
\textbf{Health culture.} Korea has a national health insurance system, so a visit to a
neighbourhood clinic (\ko{의원} uiwon) is quick and inexpensive; prescriptions are filled
at a separate \ko{약국} next door. Many people also visit practitioners of traditional
Korean medicine (\ko{한의원} hanuiwon) for acupuncture and herbal remedies. Food is
considered part of health: when you are sick you eat \ko{죽} juk (rice porridge), and
the warm advice \ko{밥 잘 챙겨 먹어} (make sure you eat properly) is a common way of
showing care. Wearing a mask when you have a cold, to protect others, was already a
common courtesy in Korea well before 2020.
\end{aside}

\section*{Review summary}
\begin{itemize}
  \item Body part + \ko{이/가 아파요}; \ko{감기에 걸렸어요}, \ko{열이 나요},
        \ko{기침을 해요}.
  \item \ko{ㅡ}-irregular: drop \ko{ㅡ}; \ko{아프다} $\to$ \ko{아파요}, \ko{크다}
        $\to$ \ko{커요}.
  \item \ko{\til (으)면}: if/when; any kind of clause can follow.
  \item \ko{\til 아/어야 하다/되다} (must), \ko{\til 아/어도 되다} (may),
        \ko{\til (으)면 안 되다} (must not), \ko{\til 지 않아도 되다} (need not),
        \ko{\til 지 마세요} (don't).
\end{itemize}
